\documentclass[aspectratio=169,10pt]{beamer}
\usetheme{default}
\setbeamertemplate{navigation symbols}{}
\usepackage{graphicx}\usepackage[T1]{fontenc}\usepackage{amsmath,amssymb}
\usepackage{listings}
\definecolor{ncsured}{HTML}{CC0000}\definecolor{ncsudk}{HTML}{990000}
\definecolor{ncsugray}{HTML}{6D6D6D}\definecolor{ncsuteal}{HTML}{008473}
\definecolor{codebg}{HTML}{F5F5F5}
\setbeamercolor{structure}{fg=ncsured}
\setbeamercolor{frametitle}{fg=white,bg=ncsured}
\setbeamercolor{title}{fg=ncsudk}
\setbeamercolor{itemize item}{fg=ncsured}
\setbeamercolor{itemize subitem}{fg=ncsured}
\setbeamerfont{frametitle}{series=\bfseries}
\setbeamertemplate{frametitle}{\nointerlineskip\vskip2pt%
  \begin{beamercolorbox}[wd=\paperwidth,ht=2.4ex,dp=1ex,leftskip=.3cm]{frametitle}%
  \small\bfseries\insertframetitle\end{beamercolorbox}}
\setbeamertemplate{footline}{\hbox{\begin{beamercolorbox}[wd=\paperwidth,ht=2.2ex,dp=1ex,leftskip=.3cm,rightskip=.3cm]{}%
  \tiny Optimal Market Making for Cryptocurrency Options \hfill Week 8 \hfill \insertframenumber\end{beamercolorbox}}}
\lstset{basicstyle=\ttfamily\scriptsize,backgroundcolor=\color{codebg},
  keywordstyle=\color{ncsured},commentstyle=\color{ncsugray},
  stringstyle=\color{ncsuteal},breaklines=true,columns=fullflexible,
  showstringspaces=false,frame=single,rulecolor=\color{ncsugray}}

\title{Week 8: Options Market Making \& Dynamic Hedging}
\subtitle{Avellaneda-Stoikov with a vector-valued (Greek) inventory}
\date{North Carolina State University \,\textbar\, Summer 2026}
\titlegraphic{
\includegraphics[width=4.2cm]{/Users/dendisuhubdy/Github/ncsu/assets/ncsu_logo.png}}
\begin{document}
\begin{frame}[plain]\titlepage\begin{center}\tiny Industry Mentor: Dendi Suhubdy (Bitwyre)  |  Guest: Fenni Kang (Coincall)\end{center}\end{frame}
\begin{frame}{Learning Objectives}
\begin{itemize}
  \item Set up the El Aoud-Abergel multi-Greek framework
  \item Understand the quadratic inventory penalty
  \item Quantify the P\&L of a delta-hedged book
\end{itemize}
\end{frame}

\begin{frame}{Lecture 1: Multi-Greek Inventory \& HJB}
\begin{itemize}
  \item Inventory becomes a vector q = (Delta, Gamma, vega, ...)
  \item Quadratic inventory penalty; reduces to AS in 1-D
  \item Multi-Greek HJB and the separation ansatz
\end{itemize}
\end{frame}

\begin{frame}{Lecture 2: Dynamic Hedging}
\begin{itemize}
  \item Continuous vs discrete delta hedging
  \item The gamma-theta-variance identity
  \item Perp futures as the hedge instrument; funding cost
\end{itemize}
\end{frame}

\begin{frame}{Illustration}
\begin{center}\includegraphics[width=0.86\textwidth]{/Users/dendisuhubdy/Github/ncsu/weekly_presentations/figures/week08.png}\end{center}
\end{frame}

\section*{Derivation}
\begin{frame}{Derivation \& Rationale}
\[\text{penalty} = \tfrac12\, q^{\mathsf T} \Sigma\, q\]

\small\textbf{\color{ncsured}Why:} Inventory risk for an options book is multi-dimensional (delta, gamma, vega...). A quadratic form is the local (mean-variance) risk of the Greek vector; $\Sigma$ is calibrated to the empirical covariance of the Greeks, so cross-risks are priced, not just diagonal ones. It reduces to AS's $q^2$ penalty when there is one Greek.
\end{frame}

\begin{frame}{Derivation \& Rationale}
\[d\Pi = \tfrac12\,\Gamma\,S^2\big(\sigma_{\text{real}}^2-\sigma_{\text{impl}}^2\big)\,dt\]

\small\textbf{\color{ncsured}Why:} The gamma-theta-variance identity: a delta-hedged option earns (or pays) the gap between realized and implied variance, scaled by dollar gamma. This single equation explains ~80% of a hedged book's daily P&L and motivates why we track gamma at all.
\end{frame}

\begin{frame}{Getting the Data from Coincall}
\begin{itemize}
  \item Greeks per contract (for the inventory vector and Sigma):
  \begin{itemize}\item \texttt{\small GET /open/option/detail/v1/\{symbol\}  -\textgreater{} delta, gamma, vega, theta, rho}\end{itemize}
  \item Hedge instrument + carry:
  \begin{itemize}\item \texttt{\small Perp book:  GET /open/futures/order/orderbook/v1/\{symbol\}}\end{itemize}
  \begin{itemize}\item \texttt{\small Funding:    GET /open/public/fundingRate/v1?symbol=BTCUSD}\end{itemize}
  \item Offline: coincall\_funding\_rates\_2025Q4.parquet for hedge-cost accounting
\end{itemize}
\vspace{2pt}\tiny\textit{Cross-check exact paths at https://docs.coincall.com/}
\end{frame}

\begin{frame}[fragile]{Code --- week8\_hedge.py}
\begin{lstlisting}
import torch
def book_greeks(positions, greeks):
    # positions: (n,), greeks: (n,5) columns [delta,gamma,vega,theta,rho]
    return (positions[:, None] * greeks).sum(0)      # aggregate inventory vector

def delta_hedge(net_delta, perp_pos, threshold=0.1):
    if net_delta.abs() > threshold:                  # rebalance only when breached
        trade = -net_delta                           # hedge with perp futures
        return perp_pos + trade, trade
    return perp_pos, torch.tensor(0.0)
\end{lstlisting}
\end{frame}
\begin{frame}{Glossary --- Key Terms}
\footnotesize\begin{description}
\item[\color{ncsured}Greek inventory vector] aggregated (delta, gamma, vega, ...) exposure of the book
\item[\color{ncsured}Quadratic penalty] the risk term (1/2) q\textasciicircum{}T Sigma q on the Greek vector
\item[\color{ncsured}Delta hedging] neutralizing directional risk using the underlying or a perpetual
\item[\color{ncsured}Gamma scalping] P\&L earned by re-hedging a long-gamma position
\item[\color{ncsured}Gamma-theta-variance identity] hedged P\&L = (1/2) Gamma S\textasciicircum{}2 (realized var - implied var) dt
\item[\color{ncsured}Funding rate] periodic payment between perpetual-swap longs and shorts
\item[\color{ncsured}El Aoud-Abergel] the multi-Greek option market-making framework
\end{description}
\end{frame}

\begin{frame}{Reading}
\begin{itemize}
  \item El Aoud-Abergel (2015); Taleb (1997); Cartea et al., Ch. 11
\end{itemize}
\end{frame}

\begin{frame}{References}
\scriptsize\begin{itemize}\setlength{\itemsep}{3pt}
  \item El Aoud, S., \& Abergel, F. (2015). A stochastic control approach to option market making. Market Microstructure and Liquidity, 1(1), 1550006.
  \item Taleb, N. N. (1997). Dynamic Hedging: Managing Vanilla and Exotic Options. Wiley.
  \item Cartea, Jaikumar, \& Penalva (2015). Algorithmic and High-Frequency Trading, Ch. 11.
\end{itemize}
\end{frame}

\end{document}